\chapter*{Homework 6}
\section*{Problem 1}
Let $X$ be a random variable such that 
$\mathbb{E}\left[\left|X\right|\right]<\infty$, 
that is $X \in L^1$. 
Denote by $\phi_X(t):= \mathbb{E}\left[e^{i t X}\right], t \in \mathbb{R}$, 
its characteristic function.
\begin{itemize}
\item[(a)]
Show that $\phi_X$ is differentiable with 
$\phi_X^{\prime}(t)=i \mathbb{E}\left[X e^{i t X}\right]$. (Hint: DCT)

\item[(b)]
If, in addition, $X$ has symmetric distribution(i.e. $X,-X$ have the same distribution),
then show that $\phi_X(t) \in \mathbb{R}$ for any $t \in \mathbb{R}$.
\end{itemize}
\begin{proof}
\begin{itemize}
\item[(a)]
Fix $t \in \mathbb{R}$. Consider the difference quotient
\[
\frac{\phi_X(t+h)-\phi_X(t)}{h}
=
\mathbb{E}\left[
e^{itX}\frac{e^{ihX}-1}{h}
\right]
\]
Notice we have (for a.e. $\omega$): 
\[
\lim_{h\to 0}\frac{e^{ihX}-1}{h}=iX
\]
Hence
\[
e^{itX}\frac{e^{ihX}-1}{h}
\longrightarrow
iXe^{itX}
\quad \text{a.s.} \quad \text{as } h\to 0
\]

Using the mean value theorem for the function
$u \mapsto e^{iuX}$, for $|h|\le 1$ we have
\[
\left|\frac{e^{ihX}-1}{h}\right|
\le |X|
\]
Also, $|e^{itX}|=1$, so
\[
\left|
e^{itX}\frac{e^{ihX}-1}{h}
\right|
=
|e^{itX}|
\left|\frac{e^{ihX}-1}{h}\right|
\le |X|
\]
Since $X \in L^1$, we have $\mathbb{E}[|X|] < \infty$, therefore  $|X|$ 
is a dominating integrable random variable for $e^{itX}\frac{e^{ihX}-1}{h}$.
Then by DCT, 
\[
\lim_{h\to 0}\frac{\phi_X(t+h)-\phi_X(t)}{h}
=
\mathbb{E}\left[\lim_{h\to 0}
e^{itX}\frac{e^{ihX}-1}{h}\right]
=
\mathbb{E}[iXe^{itX}]
\]
Thus $\phi_X$ is differentiable and
\[
\phi_X'(t)=i\mathbb{E}[Xe^{itX}]
\]

\item[(b)]
Since random
variables with the same distribution 
have the same expectation under measurable functions for which the
expectation exists, we get
\[
\phi_X(t)=\mathbb{E}[e^{itX}]=\mathbb{E}[e^{it(-X)}]=\mathbb{E}[e^{-itX}]
=\phi_X(-t)
\]
On the other hand since 
\(
\phi_X(-t)=\overline{\phi_X(t)}
\), we thus have
\[
\phi_X(t)=\overline{\phi_X(t)}
\]
A complex number equal to its own conjugate must be real. Therefore,
\[
\phi_X(t)\in \mathbb{R}, \qquad \forall t\in\mathbb{R}
\]
Writing
\[
\phi_X(t)=\mathbb{E}[\cos(tX)]+i\mathbb{E}[\sin(tX)]
\]
Since $X$ is symmetric, $\sin(tx)$ is odd, so
\[
\mathbb{E}[\sin(tX)]=0
\]
Therefore $\phi_X(t)=\mathbb{E}[\cos(tX)] \in \mathbb{R}$.
\end{itemize}
\end{proof}



\section*{Problem 2}
Let $X \sim \operatorname{Bin}(n, p)$, 
where $n \in \mathbb{N}, p \in(0,1)$ and
$Y \sim \operatorname{Pois}(\lambda)$, 
where $\lambda>0$.
\begin{itemize}
\item[(a)] Compute the characteristic functions of $X, Y$.
\item[(b)] 
Let $\left(p_n\right)_{n \in \mathbb{N}}$ be a sequence in $(0,1)$
such that $\lim _{n \rightarrow \infty} n p_n=\lambda$ 
and $\left(X_n\right)_{n \in \mathbb{N}}$ be a sequence of random variables 
 with $X_n \sim \operatorname{Bin}\left(n, p_n\right)$. 
Show that $X_n \xrightarrow{d} Y$.
\end{itemize}
\begin{proof}
\begin{itemize}
\item[(a)]
We have
\[
\mathbb{P}(X=k)=\binom{n}{k}p^k(1-p)^{n-k}, \quad k=0,1,\dots,n
\]
Therefore,
\[
\phi_X(t)
=
\mathbb{E}[e^{itX}]
=
\sum_{k=0}^n e^{itk}\binom{n}{k}p^k(1-p)^{n-k}
\]
We factor $e^{itk}$ into $(pe^{it})^k/p^k$ and obtain
\[
\phi_X(t)
=
\sum_{k=0}^n \binom{n}{k}(pe^{it})^k(1-p)^{n-k}
\]
By the binomial formula,
\[
\phi_X(t)=(1-p+pe^{it})^n
\]

Next for $Y \sim \operatorname{Pois}(\lambda)$ we have
\[
\mathbb{P}(Y=k)=e^{-\lambda}\frac{\lambda^k}{k!}, \quad k=0,1,2,\dots
\]
Hence
\[
\phi_Y(t)
=
\mathbb{E}[e^{itY}]
=
\sum_{k=0}^{\infty} e^{itk} e^{-\lambda}\frac{\lambda^k}{k!}
\]
Thus
\[
\phi_Y(t)
=
e^{-\lambda}\sum_{k=0}^{\infty}\frac{(\lambda e^{it})^k}{k!}
=
e^{-\lambda}e^{\lambda e^{it}}
=
e^{\lambda(e^{it}-1)}
\]

So the characteristic functions are
\[
\phi_X(t)=(1-p+pe^{it})^n,
\quad
\phi_Y(t)=e^{\lambda(e^{it}-1)}
\]

\item[(b)]
By part (a),
\[
\phi_{X_n}(t)=(1-p_n+p_ne^{it})^n
=
\left(1+p_n(e^{it}-1)\right)^n
\]
We now compute the limit as $n \to \infty$.

Set
\[
a_n:=p_n(e^{it}-1)
\]
Since $p_n \to 0$ (as $np_n \to \lambda < \infty$), 
we have $a_n \to 0$. Therefore,
\[
\log(1+a_n)\to a_n,
\quad n\to\infty
\]
Hence
\[
n\log(1+a_n)\to  na_n
=
np_n(e^{it}-1)
\implies
\lambda(e^{it}-1)
\]
Exponentiating, we get
\[
\phi_{X_n}(t)
=
\exp\bigl(n\log(1+a_n)\bigr)
\longrightarrow
\exp\bigl(\lambda(e^{it}-1)\bigr)
\]
But by part (a),
\[
\exp\bigl(\lambda(e^{it}-1)\bigr)=\phi_Y(t)
\]
Thus for every $t \in \mathbb{R}$,
\[
\phi_{X_n}(t)\longrightarrow \phi_Y(t)
\]

Since $\phi_Y$ is the characteristic function 
of $Y \sim \operatorname{Pois}(\lambda)$, by the uniqueness theorem
for characteristic functions, this implies
\[
X_n \xrightarrow{d} Y
\]
\end{itemize}
\end{proof}



\section*{Problem 3}
Show that

$$
\lim _{n \rightarrow \infty} \int_0^{n / 2} 
\frac{2^n}{(n-1)!} t^{n-1} e^{-2 t} dt
=\frac{1}{2}
$$
Hint: Observe that the integral is the probability of an event 
related to a Gamma distribution. Can we apply the central limit theorem?

\begin{proof}
Let
\[
I_n:=\int_0^{n/2}\frac{2^n}{(n-1)!}t^{n-1}e^{-2t}\,dt
\]
Then 
\[
f_n(t)=\frac{2^n}{(n-1)!}t^{n-1}e^{-2t}\mathbf{1}_{(0,\infty)}(t)
\]
is the density of a Gamma distribution with parameters $(n,2)$, that is,
\(
T_n\sim \Gamma(n,2)
\). 
Hence
\[
I_n=\mathbb{P}(T_n\le n/2)
\] 

Now let $X_1,X_2,\dots$ be i.i.d. random variables with
\[
X_i\sim \operatorname{Exp}(2)
\]
We know that 
\[
T_n=X_1+\cdots+X_n
\]
Also,
\[
\mu:=\mathbb{E}[X_1]=\frac12,
\quad
\sigma^2:=\operatorname{Var}(X_1)=\frac14
\]

Therefore,
\[
I_n=\mathbb{P}(X_1+\cdots+X_n\le n/2)
=
\mathbb{P}\!\left(\frac{T_n-n\mu}{\sigma\sqrt n}\le 0\right)
\]
Since $\mu=1/2$ and $\sigma=1/2$, this is
\[
I_n=\mathbb{P}\!\left(\frac{T_n-n/2}{\sqrt n/2}\le 0\right)
\]

By the Central Limit Theorem,
\[
\frac{T_n-n\mu}{\sigma\sqrt n}
=
\frac{T_n-n/2}{\sqrt n/2}
\xrightarrow{d} Z,
\quad Z\sim N(0,1)
\]
Hence, since the standard normal distribution function is continuous at $0$,
\[
\lim_{n\to\infty} I_n
=
\lim_{n\to\infty}
\mathbb{P}\!\left(\frac{T_n-n/2}{\sqrt n/2}\le 0\right)
=
\mathbb{P}(Z\le 0)
=
\frac12
\]
Thus
\[
\lim _{n \rightarrow \infty} \int_0^{n / 2}
\frac{2^n}{(n-1)!} t^{n-1} e^{-2 t} dt
=\frac{1}{2}
\]
\end{proof}





\section*{Problem 4}
A casino offers the following random game: 
A player rolls a fair die once. 
If the outcome is 2 or 4, then the player wins 3 euros from the casino. 
If the outcome is $1,3,5$, then the player loses 4 euros to the casino.
If the outcome is 6 , then the player neither wins nor loses.
If 90 players play the above game independently, 
find approximately the probability that the casino wins 
at least 30 euros in total.
\begin{solution}
Let $X_i$ be the gain of the casino from the $i$-th player, for $i=1,\dots,90$. Then the random variables
$X_1,\dots,X_{90}$ are independent and identically distributed, with
\[
X_i=
\begin{cases}
-3, & \text{if the player wins 3 euros}\\
4, & \text{if the player loses 4 euros}\\
0, & \text{if the outcome is 6}
\end{cases}
\]
Since the die is fair, we have
\[
\mathbb{P}(X_i=-3)=\frac{2}{6}=\frac13,\qquad
\mathbb{P}(X_i=4)=\frac{3}{6}=\frac12,\qquad
\mathbb{P}(X_i=0)=\frac{1}{6}
\]

Let
\[
S_{90}=X_1+\cdots+X_{90}
\]
be the total gain of the casino after 90 players. We want to approximate
\[
\mathbb{P}(S_{90}\ge 30)
\]

We first compute the mean and variance of $X_1$. The mean is
\[
\mu:=\mathbb{E}[X_1]
=(-3)\cdot \frac13+4\cdot \frac12+0\cdot \frac16
=-1+2=1
\]
Also,
\[
\mathbb{E}[X_1^2]
=9\cdot \frac13+16\cdot \frac12+0
=3+8=11
\]
Hence
\[
\sigma^2:=\operatorname{Var}(X_1)=\mathbb{E}[X_1^2]-\mu^2=11-1=10
\]

Therefore,
\[
\mathbb{E}[S_{90}]=90\mu=90,
\quad
\operatorname{Var}(S_{90})=90\sigma^2=900
\]
So the standard deviation of $S_{90}$ is
\[
\sqrt{900}=30
\]

By the Central Limit Theorem,
\[
\frac{S_{90}-90}{30}\approx N(0,1)
\]
Thus,
\[
\mathbb{P}(S_{90}\ge 30)
=
\mathbb{P}\left(\frac{S_{90}-90}{30}\ge \frac{30-90}{30}\right)
\approx
\mathbb{P}(Z\ge -2)
\]
where $Z\sim N(0,1)$. Since
\[
\mathbb{P}(Z\ge -2)=\mathbb{P}(Z\le 2)\approx 0.9772
\]
we conclude that
\[
\mathbb{P}(S_{90}\ge 30)\approx 0.9772
\]
Hence, the probability that the casino wins at least 30 euros in total is 
approximately 0.9772.
\end{solution}


\section*{Problem 5}
Assume that $\left(X_n\right)_{n \in \mathbb{N}}$ is an i.i.d. 
sequence of random variables such that $\mathbb{E}\left[X_1\right]=0$ 
and $\mathbb{E}\left[X_1^2\right]=1$. Show that
$$
\frac{\sum_{i=1}^n X_i}{\sqrt{\sum_{i=1}^n X_i^2}} 
\xrightarrow{d} Z, \quad Z \sim N(0,1) 
$$
\begin{proof}
Let
\[
S_n:=\sum_{i=1}^n X_i,
\quad
Q_n:=\sum_{i=1}^n X_i^2
\]
We want to show that
\[
\frac{S_n}{\sqrt{Q_n}} \xrightarrow{d} Z,
\quad Z\sim N(0,1)
\]

First, since $(X_n)_{n\in\mathbb{N}}$ are i.i.d. with
\(
\mathbb{E}[X_1]=0,
\) and \(
\mathbb{E}[X_1^2]=1\), 
we have $\operatorname{Var}(X_1)=1$. 
And thus by the Central Limit Theorem,
\[
\frac{S_n}{\sqrt{n}} \xrightarrow{d} Z,
\quad Z\sim N(0,1)
\]
And then, consider $Q_n$. 
Since $(X_i^2)$ are i.i.d. and $\mathbb{E}[X_1^2]=1<\infty$, 
by the Law of Large Numbers we have
\[
\frac{Q_n}{n} \xrightarrow{P} 1
\]
By continuity of the square root function,
\[
\sqrt{\frac{Q_n}{n}} \xrightarrow{P} 1
\]

Write
\[
\frac{S_n}{\sqrt{Q_n}}
=
\frac{S_n}{\sqrt{n}} \cdot \frac{1}{\sqrt{Q_n/n}}
\]

Define
\[
A_n:=\frac{S_n}{\sqrt{n}},
\quad
B_n:=\frac{1}{\sqrt{Q_n/n}}
\]
Then we have shown that
 $A_n \xrightarrow{d} Z$ and $B_n \xrightarrow{P} 1$. 


\textbf{Now we claim that: 
$A_n B_n \xrightarrow{d} Z$.}

It suffices to show that $\mathbb{E}[f(A_n B_n)] \to \mathbb{E}[f(Z)]$ 
for every bounded continuous function $f:\mathbb{R}\to\mathbb{R}$.

Let $f:\mathbb{R}\to\mathbb{R}$ be bounded and continuous.
Then 
\[
\mathbb E[f(A_nB_n)]-\mathbb E[f(Z)]
=
\bigl(\mathbb E[f(A_nB_n)]-\mathbb E[f(A_n)]\bigr)
+
\bigl(\mathbb E[f(A_n)]-\mathbb E[f(Z)]\bigr)
\]

Since $A_n\xrightarrow{d}Z$, the second term converges to $0$. 
It remains to show that
\(
\mathbb E[f(A_nB_n)]-\mathbb E[f(A_n)]\to 0
\).

Observe
\(
A_nB_n-A_n=A_n(B_n-1)
\). 
Because $A_n\xrightarrow{d}Z$, the sequence $(A_n)$ is tight. 
Also, since $B_n\xrightarrow{P}1$, we have
\[
B_n-1\xrightarrow{P}0
\]
It follows that
\[
A_n(B_n-1) = A_nB_n-A_n\xrightarrow{P}0
\]

We now show that
\[
f(A_nB_n)-f(A_n)\xrightarrow{P}0
\]
Fix $\varepsilon>0$. Since $(A_n)$ is tight, there exists $M>0$ such that
\[
\sup_{n\ge 1}\mathbb P(|A_n|>M)<\varepsilon
\]
Since $f$ is continuous on the compact interval $[-M-1,M+1]$, it is uniformly continuous there. Thus there
exists $\delta>0$ such that whenever $x,y\in[-M-1,M+1]$ and $|x-y|<\delta$, 
we have
\[
|f(x)-f(y)|<\varepsilon
\]
Now on the event
\[
\{|A_n|\le M,\ |A_nB_n-A_n|<\min(\delta,1)\}
\]
we also have $|A_nB_n|\le M+1$, so
\[
|f(A_nB_n)-f(A_n)|<\varepsilon
\]
Therefore,
\[
\mathbb P\bigl(|f(A_nB_n)-f(A_n)|>\varepsilon\bigr)
\le
\mathbb P(|A_n|>M)+\mathbb P(|A_nB_n-A_n|\ge \min(\delta,1))
\]
The first term is less than $\varepsilon$, 
and the second term tends to $0$. Hence
\[
f(A_nB_n)-f(A_n)\xrightarrow{P}0
\]

Since $f$ is bounded, the random variables $f(A_nB_n)-f(A_n)$ are uniformly 
bounded. Therefore,
\[
\mathbb E[f(A_nB_n)-f(A_n)]\to 0
\]
Combining this with
\(
\mathbb E[f(A_n)]\to \mathbb E[f(Z)]
\), 
we get
\[
\mathbb E[f(A_nB_n)]\to \mathbb E[f(Z)]
\]
Thus
\[
A_nB_n\xrightarrow{d}Z
\]
This finishes the proof that 
\[
\frac{S_n}{\sqrt{Q_n}}
=
A_nB_n
\xrightarrow{d} Z
\]
That is,
\[
\frac{\sum_{i=1}^n X_i}{\sqrt{\sum_{i=1}^n X_i^2}}
\xrightarrow{d} Z,
\quad Z\sim N(0,1)
\]
\end{proof}